\chapter{Related Work}\label{cap:related-work}
\glsresetall

% =====================================================================
% SKELETON. Target: 8--12 pages.
%
% Source material: docs/experiments/2026-07-11-related-work-study.md holds
% full per-paper analyses (method, data, metrics, headline numbers with
% section references) from complete reads of the four PDFs in related_work/.
% Pull the prose from there; do not re-derive it.
%
% Organise BY TASK FORMULATION, not chronologically. The whole point of the
% chapter is to show that the four closest systems solve three different
% problems, and that none solves the one posed in Chapter~\ref{cap:introduction}.
% =====================================================================

\todo[inline]{Opening paragraph: the three task formulations (transfer, editing, generation) and where this work sits.}

% =====================================================================
\section{Expression Transfer from an Exemplar}\label{sec:rw-transfer}

% \citet{jiang2024emojidiff}: control signal is an RGB exemplar image
% (masked face, CLIP-encoded) injected via IP-Adapter-style decoupled
% cross-attention into frozen SD1.5; identity from a reference photo;
% ArcFace identity loss on one-step decodes.
% Relation: adjacent, not competing. It cannot accept "smile = 0.5" as
% input; a user must supply an exemplar per intensity level.
% Their Sec. 1 critique of signal-conditioned methods (detail loss,
% dependence on third-party detector accuracy) applies to this work and
% should be acknowledged and answered in Cap04.

\todo[inline]{EmojiDiff: mechanism, why exemplar conditioning is not a numeric control axis.}

% =====================================================================
\section{Expression Editing of a Source Image}\label{sec:rw-editing}

\subsection{Action-Unit Conditioned Editing}\label{subsec:rw-magicface}

% \citet{wei2025magicface}: 12-dim AU *variation* vector (target minus
% source), injected through a linear layer added to the time embedding;
% FULL fine-tune of the denoising UNet plus a UNet-copy ID encoder;
% AU dropout plus CFG on the AU condition.
% Relation: closest in control semantics (continuous AU intensity on SD),
% but editing rather than generation, relative rather than absolute targets
% (requires source-AU estimation at inference), no temporal mechanism, and
% far more trainable parameters.
% Their lab-vs-wild ablation predicts this work's lab-domain artifacts.

\todo[inline]{MagicFace: relative AU vectors, full fine-tune, and what "relative" costs at inference.}

\subsection{Intensity as a Text-Embedding Direction}\label{subsec:rw-pixelsmile}

% \citet{hua2026pixelsmile}: intensity alpha scales a text-embedding
% direction on Qwen-Image-Edit with LoRA; source of the CLS protocol
% adopted in Cap04.
% Two points must be made:
%   1. Their zero-shot row (no training, text-embedding interpolation alone)
%      already achieves substantial linearity. This is the "why not just
%      prompt it?" objection and it must be answered, not ignored -- hence
%      the prompt-engineering baseline planned in Cap06.
%   2. Their negative MEAD result does NOT transfer: they mapped MEAD's
%      three discrete levels to {0.5, 0.75, 1.0}, discarding the continuous
%      per-frame signal this work extracts.

\todo[inline]{PixelSmile: the CLS protocol, the zero-shot baseline objection, why their MEAD result does not transfer.}

% =====================================================================
\section{Action-Unit Conditioned Generation}\label{sec:rw-fineface}

% THE PRIOR ART. State this plainly and early: \citet{varanka2024fineface}
% (July 2024) generates faces from text conditioned on a continuous 12-dim
% AU intensity vector, with no source image, using an adapter on a frozen
% backbone plus CFG on the AU condition, and demonstrates out-of-range
% extrapolation. Its control is BROADER than this work's single AU12 axis.
% Any claim of primacy for "AU-conditioned generation" is therefore
% unavailable and must not appear anywhere in the text.
%
% What it does not do, and this is the opening this work occupies:
%   - its intensity sweep is N INDEPENDENT stills, sharing only a seed;
%     nothing binds identity across levels architecturally;
%   - it reports NO face-identity metric at all;
%   - intensity control is shown qualitatively only, with the authors
%     conceding the scale is nonlinear; there is no commanded-vs-detected
%     correlation anywhere.

\todo[inline]{FineFace: present as prior art, precisely delimit what it leaves open.}

% =====================================================================
\section{Positioning}\label{sec:rw-positioning}

% The table below is transcribed from the comparison in
% docs/experiments/2026-07-11-related-work-study.md. Verify each cell
% against the papers before submitting, and keep the "Ours" column honest:
% the FineFace head-to-head (2026-07-21) showed ramp-following is a TIE and
% that FineFace WINS absolute calibration. The differentiator is the
% sequence-native formulation and the measurement protocol.

\begin{table}[H]
    \centering
    \caption{Positioning of this work against the closest systems.}
    \label{tab:positioning}
    \footnotesize
    \begin{tabularx}{\textwidth}{@{}l>{\raggedright\arraybackslash}X>{\raggedright\arraybackslash}X>{\raggedright\arraybackslash}X@{}}
        \toprule
        Axis & MagicFace / PixelSmile & FineFace & This work \\
        \midrule
        Task              & editing a source image      & generation, independent stills   & generation, one coherent clip \\
        Control input     & relative AU vector / unitless $\alpha$ & absolute 12-AU vector   & per-frame AU12 trajectory \\
        Source image       & required                   & not required                     & not required \\
        Identity across levels & loss term (single image) & no mechanism, no metric        & temporal attention (measured) \\
        Intensity reported as & aggregate error / VLM-judged & qualitative sweeps          & correlation + absolute error \\
        Label/eval circularity & unflagged               & unflagged                        & flagged and mitigated \\
        \bottomrule
    \end{tabularx}\\[4pt]
    {\footnotesize Source: prepared by the author.}
\end{table}

\todo[inline]{Closing paragraph: state the gap this work fills, in one sentence, and forward-reference Cap04.}
